\chapter{Convergence of Measures}

In analysis and probability one often studies \emph{objects that are not given by points but by how they distribute mass}.
A probability measure $\mu$ on $(E,\B(E))$ is exactly such an object: instead of a single outcome $x\in E$, it encodes
how likely different regions $A\subset E$ are via $\mu(A)$. We would like to give a notion of convergence to these kind of objects.

\medskip

\noindent
\textbf{Why do we need a notion of convergence?}
For instance, many limiting procedures naturally produce sequences of measures:
\begin{itemize}
\item empirical measures $\mu_n:=\frac1n\sum_{k=1}^n \delta_{X_k}$ in statistics and learning,
\item laws of approximations (Euler schemes, particle methods, discretizations),
\item distributions of random variables $X_n$ when studying weak limits (CLT, invariance principles),
\item measures as solutions to PDEs/variational problems (e.g.\ concentration effects).
\end{itemize}
To apply a limit processes in these kind of situations, we would like to make sense of statements like ``$\mu_n$ converges to $\mu$''. Therefore, we must specify \emph{what aspects of the mass distribution}
are required to converge. We can decide on that by understanding the two ways  measures compare against each other, namely
\begin{enumerate}[label=(\roman*)]
\item \emph{Setwise viewpoint:} compare probabilities of sets, i.e.\ $\mu_n(A)$ versus $\mu(A)$.
\item \emph{Test-function viewpoint:} compare integrals against a class of functions,
\[
\int f\,d\mu_n \quad \text{versus} \quad \int f\,d\mu.
\]
\end{enumerate}
The second viewpoint is often more flexible, since it interacts well with approximation arguments and functional-analytic tools.

\medskip

\noindent
\textbf{Different notions of convergence (and the intuition).}
Not all convergences are equally strong; they are tailored to different applications.
\begin{itemize}
\item \emph{Setwise convergence} requires $\mu_n(A)\to\mu(A)$ for \emph{every} measurable set $A$.
It is very strong: it forces convergence of probabilities for all events, but is typically too rigid to obtain compactness.
\item \emph{Weak convergence} (also called \emph{narrow} convergence) requires
\[
\int f\,d\mu_n \to \int f\,d\mu \qquad \forall f\in C_b(E).
\]
Intuitively, bounded continuous $f$ can only ``see'' the distribution at a coarse level:
it cannot detect changes on sets with very rough boundaries, but it captures convergence of distributions in the sense used in probability theory.
This is the natural notion behind convergence in distribution of random variables.
\item \emph{Vague convergence} is like weak convergence but only tests against $C_c(E)$.
It is useful on non-compact spaces, because compactly supported tests ignore what happens ``at infinity''.
\end{itemize}

\medskip

\noindent
\textbf{Compactness and tightness: preventing mass from escaping.}
On non-compact spaces, a sequence of probability measures may ``drift to infinity''.
Weak convergence cannot hold unless the mass stays essentially in a compact region.
The right quantitative condition is \emph{tightness}: for every $\varepsilon>0$ there exists a compact $K$ such that
$\mu_n(K)\ge 1-\varepsilon$ for all $n$.
Prokhorov-type theorems show that (on Polish spaces, and in a suitable sense on metric spaces) tightness is exactly the
compactness condition for weak convergence. Conceptually:
\[
\text{tightness} \quad \Longleftrightarrow \quad \text{``no loss of mass at infinity''} \quad
\Longrightarrow \quad \text{existence of weakly convergent subseq.}
\]

\medskip

\noindent
\textbf{Big picture.}
Convergence of measures is the language that unifies:
\emph{limits of distributions} (probability),
\emph{limits of approximations} (numerics),
and \emph{limits of mass distributions} (analysis/PDE).
The various notions differ by what test objects they allow and therefore by what phenomena they can (or cannot) detect.


\section{Setting and basic notions}

Let $(E,d)$ be a Polish space, and let $\B(E)$ be its Borel $\sigma$-algebra. 

\begin{definition}[Totally bounded and relatively compact]
Let $A\subset E$.
\begin{enumerate}[label=(\alph*)]
\item $A$ is \emph{totally bounded} if for every $\varepsilon>0$ there exist $x_1,\dots,x_N\in E$ such that
\[
A\subset \bigcup_{k=1}^N B(x_k,\varepsilon).
\]
\item $A$ is \emph{relatively compact} if $\overline{A}$ is compact.
\end{enumerate}
\end{definition}

\begin{remark}
In metric spaces one has: $A$ totally bounded $\Longleftrightarrow$ $A$ relatively compact (i.e.\ precompact).
\end{remark}

\begin{definition}[Probability measures]
We write
\[
\mathcal{M}_1(E):=\mathcal{P}(E):=\{\text{Borel probability measures on }(E,\B(E))\}.
\]
\end{definition}

\begin{proposition}[Regularity]
Let $(E,d)$ be a metric space and $\mu\in\mathcal{M}_1(E)$. Then for every Borel set $A\in\B(E)$,
\[
\mu(A)=\sup\{\mu(K): K\subset A,\ K \text{ compact}\}
=\inf\{\mu(U): A\subset U,\ U \text{ open}\}.
\]
\end{proposition}

\begin{remark}
This is the usual inner/outer regularity (Radon regularity) of Borel probability measures on metric spaces.
\end{remark}

\section{Test functions and separating classes}

% ============================================================
% Intro: Separating classes
% Insert at the beginning of Section 2 (before definitions)
% ============================================================

We are now moving toward the notion of weak convergence as we defined informally in the introduction. But before its formal definition we must define some concepts and objects needed.

A recurring theme in these notes is that probability measures are not accessed directly,
but only through the values they assign to \emph{tests}:
either to sets $A\mapsto \mu(A)$ or to functions $f\mapsto \int f\,d\mu$.
The idea of a \emph{separating class} formalizes when a restricted family of such tests
is sufficient to uniquely determine a measure.

\medskip

\noindent
\textbf{Why is this needed?}
The Borel $\sigma$-algebra $\B(E)$ is typically very large and complicated.
Requiring equality $\mu(A)=\nu(A)$ for all $A\in\B(E)$ is conceptually clear,
but not useful for defining or verifying convergence.
Instead, we want a \emph{smaller and more tractable family of tests}
that still contains enough information to recover the entire measure.

\medskip

Instead of asking for agreement on all measurable sets,
we test measures against functions:
\[
\mu \longmapsto \int f\,d\mu .
\]
If the class of functions is chosen appropriately, then knowing these integrals
for all functions in the class uniquely determines the measure.
A separating class is precisely a family of test functions with this property.

\medskip

\noindent
\textbf{Intuition behind separation.}
A class $\mathcal{C}$ separates measures if its elements can
\emph{resolve arbitrarily fine features of sets}.
In metric spaces, this is achieved by continuous functions that approximate indicators
of closed (or open) sets.
The key idea is:
\begin{quote}
although indicator functions are discontinuous and therefore excluded,
they can be approximated from above or below by continuous bounded functions.
\end{quote}
As a consequence, integrals against continuous bounded functions
encode the same information as probabilities of closed (or open) sets.

\begin{definition}[Continuous function classes]
Define
\[
C(E):=\{f:E\to\R \ :\ f \text{ continuous}\},\qquad
C_b(E):=\{f\in C(E): \|f\|_\infty<\infty\},
\]
\[
C_c(E):=\{f\in C(E): \supp(f) \text{ is compact}\},
\]
where $\|f\|_\infty:=\sup_{x\in E}|f(x)|$.
\end{definition}

\begin{definition}[Separating class]
A class $\mathcal{C}$ of bounded measurable functions $E\to\R$ is called \emph{separating} for $\mathcal{M}_1(E)$ if
\[
\Big(\forall f\in\mathcal{C}: \int f\,dP=\int f\,dQ\Big)\ \Longrightarrow\ P=Q,
\qquad P,Q\in\mathcal{M}_1(E).
\]
\end{definition}

\medskip

\noindent
\textbf{Why $C_b(E)$?}
The space $C_b(E)$ is large enough to be separating, but still manageable:
\begin{itemize}
\item boundedness ensures integrals are always well-defined,
\item continuity allows approximation arguments using the metric structure,
\item the class is stable under algebraic operations and limits.
\end{itemize}
This balance explains why weak convergence is defined using $C_b(E)$
rather than all measurable functions.

The next lemmas justify reducing the study of weak convergence of probability measures to a
small and well--structured class of test objects. They show that convergence of integrals
against bounded continuous functions is equivalent to convergence of probabilities of
closed sets, allowing one to pass seamlessly between functional and set--theoretic
characterizations of weak convergence.

The first one formalizes the idea that closed sets already contain enough information to determine a Borel probability measure on a metric space. 
Although the Borel $\sigma$-algebra is generated by all open sets, it is equally generated by closed sets, which behave better with respect to limits.
If two measures agree on all closed sets, then they agree on a generating class that is stable under finite intersections.
Measure-theoretic structure theorems (specifically the $\pi$--$\lambda$ theorem) then guarantee that this agreement propagates to the entire Borel $\sigma$-algebra.
Conceptually, closed sets act as ``building blocks'' of measurable sets, so knowing the mass assigned to every closed set uniquely fixes the whole measure.



\begin{lemma}[Closed sets determine a probability measure]
\label{lem:closed_det}
Let $P,Q\in\mathcal{M}_1(E)$. If
\[
P(F)=Q(F)\quad\text{for all closed }F\subset E,
\]
then $P=Q$.
\end{lemma}

\begin{proof}
Let
\[
\mathcal{D}:=\{A\in\B(E): P(A)=Q(A)\}.
\]
Then $\mathcal{D}$ is a Dynkin system ($\lambda$-system): $E\in\mathcal{D}$, and if $A\in\mathcal{D}$ then $A^c\in\mathcal{D}$, and if $(A_n)$ are disjoint in $\mathcal{D}$ then $\cup_n A_n\in\mathcal{D}$.

Let $\mathcal{C}$ be the family of closed sets. Then $\mathcal{C}$ is a $\pi$-system since the intersection of two closed sets is closed. By assumption, $\mathcal{C}\subset\mathcal{D}$.

Finally, $\sigma(\mathcal{C})=\B(E)$ because closed sets generate the Borel $\sigma$-algebra. Hence by the $\pi$--$\lambda$ theorem, $\mathcal{D}$ contains $\B(E)$, so $P(A)=Q(A)$ for all $A\in\B(E)$, i.e.\ $P=Q$.
\end{proof}

The second lemma explains why testing two probability measures against all bounded continuous functions is strong enough to recover their values on closed sets.
Although the indicator function $\mathbf{1}_F$ of a closed set $F$ is discontinuous and therefore cannot be used directly as a test function,
it can be approximated from above by a sequence of bounded continuous functions that ``blur'' the boundary of $F$ at smaller and smaller scales.
These approximations equal $1$ on $F$, decay continuously to $0$ in a thin neighborhood outside $F$, and vanish away from $F$.
If two measures agree on all bounded continuous test functions, then they agree on each approximating function, and by monotone convergence this forces them to agree on the limit, namely on the indicator of $F$ itself.
Thus, even though closed sets are not tested directly, their probabilities are fully encoded in the integrals against $C_b(E)$.


\begin{lemma}[From $C_b(E)$ to closed sets ]
\label{lem:approx_closed}
Let $P,Q\in\mathcal{M}_1(E)$. If
\[
\int_E f\,dP=\int_E f\,dQ\quad\text{for all }f\in C_b(E),
\]
then $P(F)=Q(F)$ for every closed $F\subset E$.
\end{lemma}

\begin{proof}
Fix closed $F\subset E$. Define $d(x,F):=\inf_{y\in F} d(x,y)$, which is continuous on $E$.
For $n\in\mathbb{N}$ define
\[
f_n(x):=\begin{cases}
1, & x\in F,\\[2pt]
1-n\,d(x,F), & 0<d(x,F)\le \frac1n,\\[2pt]
0, & d(x,F)>\frac1n.
\end{cases}
\]
Because $x\mapsto d(x,F)$ is continuous and the piecewise formula glues continuously at the boundaries,
we have $f_n\in C_b(E)$ and $0\le f_n\le 1$.

Moreover $f_n(x)\uparrow 1_{F}(x)$ pointwise:
if $x\in F$ then $f_n(x)=1$ for all $n$; if $x\notin F$ then $d(x,F)=:r>0$ and for $n>1/r$ one has $d(x,F)>1/n$, hence $f_n(x)=0$.

By assumption, $\int f_n\,dP=\int f_n\,dQ$ for all $n$. By monotone convergence,
\[
P(F)=\int 1_{F}\,dP=\lim_{n\to\infty}\int f_n\,dP
=\lim_{n\to\infty}\int f_n\,dQ=\int 1_{F}\,dQ=Q(F).
\]
\end{proof}

\begin{corollary}[$C_b(E)$ is separating for $\mathcal{M}_1(E)$]
If $\int f\,dP=\int f\,dQ$ for all $f\in C_b(E)$, then $P=Q$.
\end{corollary}

\begin{proof}
Lemma~\ref{lem:approx_closed} gives $P(F)=Q(F)$ for all closed $F$, then Lemma~\ref{lem:closed_det} gives $P=Q$.
\end{proof}

Separating classes are the conceptual foundation of next topics, i.e., weak convergence,
compactness results, and Prokhorov-type theorems.
They explain why convergence of integrals against a restricted family of functions
is strong enough to identify limits,
and they justify replacing complicated setwise arguments by functional ones.

\section{Weak and vague convergence}

In many applications, requiring convergence of measures on all measurable sets is too strong and unstable under limits.
Weak convergence relaxes this by testing measures only against bounded continuous functions, which probe the distribution at a macroscopic level and ignore small irregularities of set boundaries.
This notion is strong enough to uniquely identify the limit measure, thanks to the separating property of $C_b(E)$, yet weak enough to enjoy compactness properties such as Prokhorov’s theorem (next section).
Vague convergence further weakens the test class to compactly supported continuous functions, allowing mass to ``escape to infinity'' and making it useful on non-compact spaces.
Both notions formalize the idea that convergence of distributions should be understood through convergence of expectations of suitable observables.

\begin{definition}[Weak convergence]
A sequence $(P_n)\subset\mathcal{M}_1(E)$ converges \emph{weakly} to $P\in\mathcal{M}_1(E)$, written $P_n\Rightarrow P$, if
\[
\int_E f\,dP_n \longrightarrow \int_E f\,dP
\qquad\text{for all } f\in C_b(E).
\]
\end{definition}

\begin{definition}[Vague convergence]
A sequence $(P_n)\subset\mathcal{M}_1(E)$ converges \emph{vaguely} to $P$ if
\[
\int_E f\,dP_n \longrightarrow \int_E f\,dP
\qquad\text{for all } f\in C_c(E).
\]
\end{definition}

\begin{remark}
The reason $C_b(E)$ appears in the definition is precisely that it is a separating class for $\mathcal{M}_1(E)$.
\end{remark}

\section{The Prokhorov metric}

Weak convergence is defined through convergence of integrals against test functions, which makes it well suited for analysis but rather indirect from a geometric point of view.
The Prokhorov metric provides a way to \emph{quantify} weak convergence by measuring how close two probability measures are in terms of the mass they assign to sets.
Its role is to turn the topology of weak convergence into a metric structure, allowing one to speak of Cauchy sequences, completeness, and separability in $\mathcal{P}(X)$.

\medskip

\noindent
The guiding idea is the following: two probability measures should be close if every closed set that has large mass under one measure also has almost as large mass under the other, up to a small error and a small spatial perturbation ($\varepsilon$-enlargement).
This perturbation is necessary because weak convergence is insensitive to small changes near boundaries, and exact set inclusion is too rigid to capture this behavior.

\medskip

\noindent
\textbf{Intuition behind the small spatial perturbation ($\varepsilon$-enlargement).}
Given a set $A\subset X$, its $\varepsilon$-enlargement $A^\varepsilon$ consists of all points whose distance to $A$ is less than $\varepsilon$.
Geometrically, $A^\varepsilon$ is obtained by ``thickening'' $A$ by a layer of radius $\varepsilon$.
In the definition of the Prokhorov metric, allowing $F$ to be replaced by $F^\varepsilon$ reflects the fact that weak convergence only controls measures up to small spatial errors near the boundary of sets.
Thus, the $\varepsilon$-enlargement compensates for boundary effects, while the additive $\varepsilon$ term controls the remaining discrepancy in mass.
Together, these two tolerances capture precisely the flexibility inherent in weak convergence.


\begin{definition}[$\varepsilon$-enlargement]
For $A\subset E$ and $\varepsilon>0$ define
\[
A^\varepsilon:=\{x\in E:\ d(x,A)<\varepsilon\},
\qquad d(x,A):=\inf_{y\in A} d(x,y).
\]
If $F$ is closed, then $F^\varepsilon$ is open.
\end{definition}

\begin{definition}[Prokhorov metric]
For $P,Q\in\mathcal{M}_1(E)$ define
\[
\beta(P,Q):=\inf\Big\{\varepsilon>0:\ P(F)\le Q(F^\varepsilon)+\varepsilon\ \ \forall \text{ closed }F\subset E\Big\}.
\]
\end{definition}

We will show that $\beta(P,Q)$ is a metric, for that we will need the following lemma.
\begin{lemma}[Key implication used for symmetry]
\label{lem:key_sym}
Let $P,Q\in\mathcal{M}_1(E)$ and let $\eta>0$. Assume that
\[
P(F)\le Q(F^\eta)+\eta\qquad\forall\text{ closed }F\subset E.
\]
Then also
\[
Q(F)\le P(F^\eta)+\eta\qquad\forall\text{ closed }F\subset E.
\]
\end{lemma}

\begin{proof}
Fix closed $F$. Then $E\setminus F^\eta$ is closed. Apply the assumption to $E\setminus F^\eta$:
\[
P(E\setminus F^\eta)\le Q\big((E\setminus F^\eta)^\eta\big)+\eta.
\]
We claim $(E\setminus F^\eta)^\eta\subset E\setminus F$. Indeed, if $x\in F$ then $d(x,E\setminus F^\eta)\ge \eta$
(because $B(x,\eta)\subset F^\eta$), hence $x\notin (E\setminus F^\eta)^\eta$. Thus
\[
Q\big((E\setminus F^\eta)^\eta\big)\le Q(E\setminus F).
\]
So
\[
P(E\setminus F^\eta)\le Q(E\setminus F)+\eta.
\]
Using $P(E)=Q(E)=1$ and rearranging gives $Q(F)\le P(F^\eta)+\eta$.
\end{proof}

\begin{proposition}[$\beta$ is a metric on $\mathcal{M}_1(E)$]
\label{prop:beta_metric}
The function $\beta$ satisfies:
\begin{enumerate}[label=(\alph*)]
\item $\beta(P,Q)\ge 0$ and $\beta(P,Q)=\beta(Q,P)$.
\item If $\beta(P,Q)=0$, then $P=Q$.
\item $\beta(P,R)\le \beta(P,Q)+\beta(Q,R)$ (triangle inequality).
\end{enumerate}
\end{proposition}

\begin{proof}
(a) Nonnegativity is immediate. For symmetry, if $\varepsilon$ satisfies
$P(F)\le Q(F^\varepsilon)+\varepsilon$ for all closed $F$, then Lemma~\ref{lem:key_sym} gives the reverse inequality,
so the same $\varepsilon$ works for $(Q,P)$; hence $\beta(P,Q)=\beta(Q,P)$.

(b) If $\beta(P,Q)=0$, then for every $\varepsilon>0$,
\[
P(F)\le Q(F^\varepsilon)+\varepsilon\qquad\forall\text{ closed }F.
\]
Fix closed $F$ and let $\varepsilon_m=1/m$. Since $F^{1/m}\downarrow F$ and measures are continuous from above,
\[
Q(F)=\lim_{m\to\infty} Q(F^{1/m}).
\]
Thus, for each $m$,
\[
P(F)\le Q(F^{1/m})+\frac1m,
\]
and letting $m\to\infty$ gives $P(F)\le Q(F)$. By symmetry (part (a)) we also get $Q(F)\le P(F)$, hence $P(F)=Q(F)$
for all closed $F$. Lemma~\ref{lem:closed_det} yields $P=Q$.

(c) Let $\alpha>\beta(P,Q)$ and $\gamma>\beta(Q,R)$. Then for every closed $F$,
\[
P(F)\le Q(F^\alpha)+\alpha \le R\big((F^\alpha)^\gamma\big)+\gamma+\alpha.
\]
But $(F^\alpha)^\gamma\subset F^{\alpha+\gamma}$, hence
\[
P(F)\le R(F^{\alpha+\gamma})+(\alpha+\gamma)\qquad\forall\text{ closed }F.
\]
Therefore $\beta(P,R)\le \alpha+\gamma$. Letting $\alpha\downarrow\beta(P,Q)$ and $\gamma\downarrow\beta(Q,R)$ gives the triangle inequality.
\end{proof}

The next theorem (stated without a formal proof) clarifies the geometric and topological structure of the space of probability measures equipped with the Prokhorov metric.
Once weak convergence has been metrized by $\beta$, it becomes essential to know whether the resulting metric space is well behaved:
separability ensures that convergence can be tested on countable families and enables diagonal arguments,
while completeness guarantees that Cauchy sequences of probability measures admit limits.
Without these properties, compactness and subsequence arguments central to Prokhorov’s theorem (next section) would fail at a foundational level.
Thus, Theorem~\ref{thm:underlyng_space_measure_space_rel} shows that when the underlying space $E$ is Polish, the space $\mathcal{M}_1(E)$ equipped with $\beta$
inherits the same favorable structure, making it a natural and robust setting for studying weak convergence and tightness.


\begin{theorem}[Relation $E$ to $\mathcal{M}_1(E)$]\label{thm:underlyng_space_measure_space_rel}
If $E$ is separable, then $(\mathcal{M}_1(E),\beta)$ is separable. If $E$ is complete, then $(\mathcal{M}_1(E),\beta)$ is complete.
Hence if $E$ is Polish, then $(\mathcal{M}_1(E),\beta)$ is Polish.
\end{theorem}

\begin{proof}
These are standard results about the Prokhorov metric; we include a clean proof sketch.

If $E$ is separable, choose a countable dense set $D$ and consider finitely supported probability measures supported on $D$ with rational weights.
This set is countable; one proves it is dense in $(\mathcal{M}_1(E),\beta)$.

If $E$ is complete, every $\beta$-Cauchy sequence is tight and admits a weak limit; one then checks convergence in $\beta$ to that limit.
Completeness of $E$ ensures the limit is a Borel probability on $E$ (no mass “escapes” to a completion).
\end{proof}

\section{Portmanteau theorem}

The definition of weak convergence is formulated through test
functions, and in applications it is often more convenient to verify convergence through sets or to
work with an explicit metric.  The Portmanteau theorem below provides exactly this bridge: it gives a
collection of equivalent criteria for $\mathbb{P}_n\Rightarrow\mathbb{P}$, ranging from convergence
of integrals against (uniformly) continuous bounded functions to one-sided bounds on probabilities
of closed/open sets and convergence on $\mathbb{P}$-continuity sets.  The additional link to the
Prokhorov metric is particularly useful because it turns weak convergence into a quantitative,
metric statement; this enables compactness arguments (via tightness/Prokhorov) and stability
estimates when passing to limits in probabilistic models.

\begin{theorem}[Portmanteau theorem and Prokhorov metric]\label{thm:portmanteau-prokhorov}
Let $(E,d)$ be a metric space and let $\mathcal{E}:=\mathcal{B}(E)$ be its Borel $\sigma$-algebra.
Let $(\mathbb{P}_n)_{n\in\mathbb{N}}\subset \mathcal{P}(E)$ and $\mathbb{P}\in\mathcal{P}(E)$ be Borel probability measures.
Consider the following statements:
\begin{enumerate}[label=\textnormal{(\alph*)}, leftmargin=3.2em]
  \item $\displaystyle \lim_{n\to\infty}\beta(\mathbb{P}_n,\mathbb{P})=0$, where $\beta$ denotes the Prokhorov metric.
  \item $\mathbb{P}_n \Rightarrow \mathbb{P}$ (weak convergence).
  \item $\displaystyle \lim_{n\to\infty}\int_E f\,d\mathbb{P}_n=\int_E f\,d\mathbb{P}\qquad \forall\, f\in C_b(E)$,
        where $f$ is uniformily continuous.
  \item $\displaystyle \limsup_{n\to\infty}\mathbb{P}_n(F)\le \mathbb{P}(F)\qquad \forall\, F\in\mathcal{E}\ \text{closed}$.
  \item $\displaystyle \liminf_{n\to\infty}\mathbb{P}_n(G)\ge \mathbb{P}(G)\qquad \forall\, G\in\mathcal{E}\ \text{open}$.
  \item $\displaystyle \lim_{n\to\infty}\mathbb{P}_n(A)=\mathbb{P}(A)\qquad \forall\, A\in\mathcal{E}\ \text{with }\mathbb{P}(\partial A)=0$,
        i.e.\ for all $\mathbb{P}$-continuity sets $A$.
\end{enumerate}
Then \textnormal{(b)}--\textnormal{(f)} are equivalent, and \textnormal{(a)} implies \textnormal{(b)}.
Moreover, if $(E,d)$ is separable, then \textnormal{(a)}--\textnormal{(f)} are equivalent (in particular, $\beta$ metrizes weak convergence).
\end{theorem}

\begin{proof}[Equivalence of \textnormal{(b)}--\textnormal{(f)}]
We prove the chain of implications
\[
\textnormal{(b)}\Rightarrow\textnormal{(c)}\Rightarrow\textnormal{(d)}\Rightarrow\textnormal{(e)}
\Rightarrow\textnormal{(f)}\Rightarrow\textnormal{(b)},
\]

\medskip
\LabelProof{\textnormal{(b)} $\Rightarrow$ \textnormal{(c)}}
If $\mathbb{P}_n\Rightarrow \mathbb{P}$, then by definition (weak convergence on a metric space)
\[
\int_E f\,d\mathbb{P}_n \to \int_E f\,d\mathbb{P}\qquad\forall f\in C_b(E).
\]
In particular, this holds for every bounded \emph{uniformly} continuous $f$, hence \textnormal{(c)} follows. 

\medskip
\LabelProof{\textnormal{(c)} $\Rightarrow$ \textnormal{(d)}}
Fix a closed set $F\subset E$ and $\delta>0$. Define the open $\delta$-neighbourhood
\[
G_\delta := F^\delta := \{y\in E:\ \exists x\in F\ \text{with}\ d(x,y)<\delta\},
\]
and define the function
\[
f_{F,\delta}(x):=\Bigl(1-\frac{d(x,F)}{\delta}\Bigr)_+,
\qquad d(x,F):=\inf_{y\in F} d(x,y).
\]
Then $f_{F,\delta}$ is bounded and (in fact Lipschitz, hence uniformly continuous), satisfies
$f_{F,\delta}=1$ on $F$ and $f_{F,\delta}=0$ on $G_\delta^{\,c}$, and $0\le f_{F,\delta}\le 1$.
Therefore
\[
\mathbb{P}_n(F)\le \int_E f_{F,\delta}\,d\mathbb{P}_n.
\]
Taking $\limsup_{n\to\infty}$ and using \textnormal{(c)} gives
\[
\limsup_{n\to\infty}\mathbb{P}_n(F)
\le \int_E f_{F,\delta}\,d\mathbb{P}
\le \mathbb{P}(G_\delta).
\]
Since $G_\delta\downarrow F$ as $\delta\downarrow 0$, continuity from above yields
$\mathbb{P}(G_\delta)\downarrow \mathbb{P}(F)$, hence letting $\delta\downarrow 0$ gives
\[
\limsup_{n\to\infty}\mathbb{P}_n(F)\le \mathbb{P}(F),
\]
which is \textnormal{(d)}. 

\medskip
\LabelProof{\textnormal{(d)} $\Rightarrow$ \textnormal{(e)}}
If $G$ is open, then $G^c$ is closed, so by \textnormal{(d)},
\[
\limsup_{n\to\infty}\mathbb{P}_n(G^c)\le \mathbb{P}(G^c).
\]
Using $\mathbb{P}_n(G)=1-\mathbb{P}_n(G^c)$ and $\mathbb{P}(G)=1-\mathbb{P}(G^c)$ yields
\[
\liminf_{n\to\infty}\mathbb{P}_n(G)\ge \mathbb{P}(G),
\]
i.e.\ \textnormal{(e)}. 

\medskip
\LabelProof{\textnormal{(d)} and \textnormal{(e)} $\Rightarrow$ \textnormal{(f)}}
Let $A\in\mathcal{E}$ satisfy $\mathbb{P}(\partial A)=0$. Then $A^\circ\subset A\subset \overline{A}$
and $\overline{A}\setminus A^\circ\subset \partial A$, hence
$\mathbb{P}(\overline{A})=\mathbb{P}(A^\circ)=\mathbb{P}(A)$.
Applying \textnormal{(d)} to $\overline{A}$ and \textnormal{(e)} to $A^\circ$ gives
\[
\mathbb{P}(A)=\mathbb{P}(\overline{A}) \ge \limsup_{n\to\infty}\mathbb{P}_n(\overline{A})
\ge \limsup_{n\to\infty}\mathbb{P}_n(A),
\]
and
\[
\liminf_{n\to\infty}\mathbb{P}_n(A)\ge \liminf_{n\to\infty}\mathbb{P}_n(A^\circ)
\ge \mathbb{P}(A^\circ)=\mathbb{P}(A).
\]
Thus $\mathbb{P}_n(A)\to \mathbb{P}(A)$, i.e.\ \textnormal{(f)} holds. 

\medskip
\LabelProof{\textnormal{(f)} $\Rightarrow$ \textnormal{(b)}.}
We show that \textnormal{(f)} implies convergence of integrals for all $f\in C_b(E)$, which is
equivalent to \textnormal{(b)} by Definition~5.10 in the notes. 
Let $f\in C_b(E)$ and write $f=f^+-f^-$, so it suffices to assume $f\ge 0$.
Set $C:=\|f\|_\infty$. Using the standard identity for nonnegative bounded functions,
\[
\int_E f\,d\mathbb{P}=\int_{0}^{C}\mathbb{P}(f>t)\,dt,
\qquad
\int_E f\,d\mathbb{P}_n=\int_{0}^{C}\mathbb{P}_n(f>t)\,dt.
\]
For each $t\in(0,C)$ one has
\[
\partial\{f>t\}\subset \{f=t\}.
\]
The level sets $\{f=t\}$ are pairwise disjoint in $t$, hence there are at most countably many
$t$ with $\mathbb{P}(f=t)>0$. Therefore, for Lebesgue-a.e.\ $t\in(0,C)$,
\[
\mathbb{P}(\partial\{f>t\})=0,
\]
so $\{f>t\}$ is a $\mathbb{P}$-continuity set. By \textnormal{(f)} we then have
$\mathbb{P}_n(f>t)\to \mathbb{P}(f>t)$ for a.e.\ $t$.
Since $0\le \mathbb{P}_n(f>t)\le 1$, dominated convergence yields
\[
\int_E f\,d\mathbb{P}_n
=\int_{0}^{C}\mathbb{P}_n(f>t)\,dt
\longrightarrow
\int_{0}^{C}\mathbb{P}(f>t)\,dt
=\int_E f\,d\mathbb{P}.
\]
Thus $\int f\,d\mathbb{P}_n\to \int f\,d\mathbb{P}$ for all $f\in C_b(E)$, i.e.\ $\mathbb{P}_n\Rightarrow \mathbb{P}$.
This proves \textnormal{(f)} $\Rightarrow$ \textnormal{(b)}. 

\medskip
Combining Steps~1--5 gives the equivalence of \textnormal{(b)}--\textnormal{(f)}.
\end{proof}




\begin{proof}[Proof of \textnormal{(a)} $\Rightarrow$ \textnormal{(b)}]
Assume \textnormal{(a)}, i.e.\ $\beta(\mathbb{P}_n,\mathbb{P})\to 0$ as $n\to\infty$, where $\beta$ is the Prokhorov metric.
Fix a closed set $F\in\mathcal{E}$. By the definition of $\beta$, for each $n$ there exists $\varepsilon_n\downarrow 0$
(with $\varepsilon_n \ge \beta(\mathbb{P}_n,\mathbb{P})$) such that
\[
\mathbb{P}_n(F)\le \mathbb{P}(F^{\varepsilon_n})+\varepsilon_n,
\]
where $F^{\varepsilon}:=\{x\in E:\ d(x,F)<\varepsilon\}$ denotes the open $\varepsilon$-neighbourhood of $F$.
Taking $\limsup_{n\to\infty}$ yields
\[
\limsup_{n\to\infty}\mathbb{P}_n(F)
\le \limsup_{n\to\infty}\mathbb{P}(F^{\varepsilon_n})+\varepsilon_n.
\]
Since $\varepsilon_n\downarrow 0$, we have $F^{\varepsilon_n}\downarrow F$ and hence, by continuity from above of
the probability measure $\mathbb{P}$,
\[
\lim_{n\to\infty}\mathbb{P}(F^{\varepsilon_n})=\mathbb{P}(F).
\]
Therefore $\limsup_{n\to\infty}\mathbb{P}_n(F)\le \mathbb{P}(F)$ for every closed $F$.
By the $(d)$, this implies $\mathbb{P}_n \Rightarrow \mathbb{P}$,
i.e.\ \textnormal{(b)} holds.
\end{proof}


\begin{proof}[Proof of \textnormal{(b)} $\Rightarrow$ \textnormal{(a)}]
Assume \textnormal{(b)}, i.e.\ $\mathbb{P}_n \Rightarrow \mathbb{P}$.
Assume moreover that $(E,d)$ is separable.

Fix $\varepsilon>0$. By separability, there exists a countable collection of open balls
$\{A_k^\varepsilon\}_{k\in\mathbb{N}}$ with diameter $<\varepsilon$ such that $\bigcup_{k=1}^\infty A_k^\varepsilon = E$.
Choose $N_1(\varepsilon)\in\mathbb{N}$ so that
\[
\mathbb{P}\Big(\bigcup_{k=1}^{N_1(\varepsilon)} A_k^\varepsilon\Big) > 1-\varepsilon,
\qquad\text{equivalently}\qquad
\mathbb{P}\Big(\bigcup_{k>N_1(\varepsilon)} A_k^\varepsilon\Big) < \varepsilon.
\]
Next, by $(e)$, for any open set $G$ there exists
$N_2(\varepsilon,G)$ such that for all $n\ge N_2(\varepsilon,G)$,
\[
\mathbb{P}_n(G) \ge \mathbb{P}(G)-\varepsilon.
\]
Using the countable basis construction, one obtains an index $N(\varepsilon)$ such that
for all $n\ge N(\varepsilon)$ and all Borel sets $A\in\mathcal{E}$,
\begin{equation}\label{eq:prokhorov_one_side}
\mathbb{P}(A)\le \mathbb{P}_n(A^{2\varepsilon}) + 2\varepsilon.
\end{equation}

\LabelProof{***Briefly: cover $A$ by unions of the finitely many balls $A_1^\varepsilon,\dots,A_{N_1(\varepsilon)}^\varepsilon$
that intersect $A$, plus the tail $\bigcup_{k>N_1(\varepsilon)}A_k^\varepsilon$ of $\mathbb{P}$-mass $<\varepsilon$;
call this union $G_0$. Then $G_0\subset A^\varepsilon$ and $\mathbb{P}(A)\le \mathbb{P}(G_0)+\varepsilon
\le \mathbb{P}_n(G_0)+2\varepsilon \le \mathbb{P}_n(A^{2\varepsilon})+2\varepsilon$.***}
\medskip

From \eqref{eq:prokhorov_one_side} we derive the reverse inequality as follows: for any Borel set $B$,
\[
\mathbb{P}(B^{\varepsilon})
= 1-\mathbb{P}\big((B^{\varepsilon})^c\big)
\ge 1-\mathbb{P}_n\big(((B^{\varepsilon})^c)^{2\varepsilon}\big)-2\varepsilon
\ge \mathbb{P}_n(B)-2\varepsilon,
\]
where we used \eqref{eq:prokhorov_one_side} with $A=(B^\varepsilon)^c$ and the set inclusion
$B \subset \big(((B^{\varepsilon})^c)^{2\varepsilon}\big)^c$ (equivalently,
$\big(((B^{\varepsilon})^c)^{2\varepsilon}\big) \subset B^c$).
Hence, for all $n\ge N(\varepsilon)$ and all Borel $B$,
\[
\mathbb{P}_n(B)\le \mathbb{P}(B^{\varepsilon})+2\varepsilon.
\]
Together with \eqref{eq:prokhorov_one_side} (renaming $A$ as $B$) this shows that for all $n\ge N(\varepsilon)$,
\[
\beta(\mathbb{P}_n,\mathbb{P}) \le 2\varepsilon.
\]
Since $\varepsilon>0$ was arbitrary, we conclude $\beta(\mathbb{P}_n,\mathbb{P})\to 0$, i.e.\ \textnormal{(a)} holds.
\end{proof}

\chapter{Tightness}

On non-compact spaces, weak convergence alone does not prevent probability mass from drifting off to infinity.
A sequence of probability measures may converge weakly on every bounded region, yet lose mass globally, so that no limit probability measure exists.
Tightness is the precise condition that rules out this phenomenon: it requires that, uniformly in the sequence, almost all mass is concentrated in a common compact set.
In this sense, tightness is the probabilistic analogue of precompactness.

\begin{definition}[Tightness]\label{def:tightness}
A family $\mathcal{K}\subset\mathcal{M}_1(E)$ is \emph{tight} if for every $\varepsilon>0$ there exists a compact set $K\subset E$ such that
\[
\inf_{P\in\mathcal{K}} P(K) > 1-\varepsilon.
\]
\end{definition}

\begin{lemma}[Polish spaces: every probability measure is tight]
Let $E$ be Polish. Then every $P\in\mathcal{M}_1(E)$ is tight. Consequently, every finite subset of $\mathcal{M}_1(E)$ is tight.
\end{lemma}

\begin{proof}
On Polish spaces, Borel probability measures are Radon, hence tight: for each $\varepsilon>0$ there exists compact $K$ with $P(K)>1-\varepsilon$.
For a finite set $\{P_1,\dots,P_m\}$, choose compacts $K_i$ with $P_i(K_i)>1-\varepsilon$, and take $K:=\cap_{i=1}^m K_i$.
Then $K$ is compact and $P_i(K)\ge 1-\varepsilon$ for all $i$.
\end{proof}

\begin{remark}[Alternative definition of tightness]
    Another reading of tightness is: for every tolerance $\epsilon > 0$, there exists a single compact set $K$ such that all measures in the family place \emph{at least} $1-\epsilon$ of their mass inside $K$, i.e., \emph{all} measures in the family have its mass inside $K$ except for a (small) $\epsilon$. Therefore, we can think in the other way around and define tightness on the part which we left out using $\sup_{P\in\mathcal{K}}P(K^c) < \epsilon$.
\end{remark}

\medskip

Prokhorov has identified tightness as the exact compactness criterion for weak convergence.
On Polish spaces, his theorem (below) states that a family of probability measures is relatively compact for the weak topology if and only if it is tight.
His theorem thus provides the missing link between analytic convergence (defined via test functions) and geometric control of mass (via compact sets).

\medskip


Conceptually, the theorem works by separating two difficulties:
first, tightness ensures that mass cannot escape to infinity, so that potential limit points are genuine probability measures;
second, once tightness is available, weak compactness follows from the metrizability and completeness properties established earlier via the Prokhorov metric.
In this way, Prokhorov’s theorem explains why weak convergence behaves like ordinary compactness in finite-dimensional settings, despite the infinite-dimensional nature of $\mathcal{P}(X)$.




\paragraph{Intuition behind the proof of Prokhorov’s theorem.}
The proof of Prokhorov’s theorem hinges on the idea that tightness provides exactly the compactness needed to extract convergent subsequences.
Tightness guarantees that, uniformly in the family, almost all mass lies in a sequence of compact sets, preventing escape to infinity.
On each such compact region, probability measures behave much like measures on a compact space, where weak compactness is easier to obtain.
Using separability, one tests convergence only on a countable convergence-determining class of functions and applies a diagonal argument to extract a subsequence whose integrals converge simultaneously for all tests.
The limiting linear functional then corresponds to a probability measure, and tightness ensures that no mass is lost in the limit, completing the argument.

\begin{theorem}[Prokhorov's theorem]
\label{thm:prokhorov}
Let $E$ be Polish and $\mathcal{K}\subset\mathcal{M}_1(E)$. The following are equivalent:
\begin{enumerate}[label=(\alph*)]
\item $\mathcal{K}$ is tight.
\item For every $\varepsilon>0$ there exists compact $K\subset E$ with $\inf_{P\in\mathcal{K}}P(K)>1-\varepsilon$.
\item $\mathcal{K}$ is relatively compact for weak convergence.
\end{enumerate}
\end{theorem}

\begin{proof}
\LabelProof{(a)$\Leftrightarrow$(b)} It is just the definition.

\LabelProof{(b)$\Rightarrow$(c)} Let $(P_n)\subset\mathcal{K}$. Tightness provides compacts sets $(K_m)_{m \geq 1}$ such that
$\inf_n P_n(K_m)\ge 1-2^{-m}$. Since $E$ is separable, there exists a countable convergence-determining class
$\mathcal{C}\subset C_b(E)$. Next, $\int f dP_n$ is bounded, thus by Bolzano-Weierstrass, there exists a subsequence $(P_{n_k})$ such that $\int f\,dP_{n_k}$
converges for all $f\in\mathcal{C}$. We know that, if $(\mu_n) \in \mathcal{M}_1(E)$ is tight and $\int f d\mu_n$ converges for all $f$ in a countable separting class, then there exist $\mu \in \mathcal{M}_1(E)$ such that $\int f d\mu_n \to \int f d\mu$ for all $f\in\mathcal{C}$.
This identifies a unique limit probability measure $P$ (because $C_b(E)$ separates measures),
and one concludes $P_{n_k}\Rightarrow P$.

\LabelProof{(c)$\Rightarrow$(a)} If $\mathcal{K}$ were not tight, there would exist $\varepsilon_0>0$ such that for every compact $K$
some $P\in\mathcal{K}$ satisfies $P(K)\le 1-\varepsilon_0$. Choose an increasing compact exhaustion $(K_m)$ of $E$
and pick $P_m\in\mathcal{K}$ with $P_m(K_m)\le 1-\varepsilon_0$.
Relative compactness yields a weakly convergent subsequence $P_{m_j}\Rightarrow P$.
Using the Portmanteau inequality for closed sets (applied to each fixed $K_m$) one deduces $P(K_m)\ge 1-\varepsilon_0$ for all $m$,
hence $P(\cup_m K_m)\ge 1-\varepsilon_0$. But $\cup_m K_m=E$, so $P(E)\ge 1-\varepsilon_0$, contradicting $P(E)=1$ unless $\varepsilon_0=0$.
Thus $\mathcal{K}$ must be tight.
\end{proof}

\begin{theorem}[Prokhorov theorem on metric spaces]
\label{thm:prokhorov_metric}
Let $(X,d)$ be a metric space and let $\mathcal{F}\subset \mathcal{P}(X)$.
If $\mathcal{F}$ is tight, then $\mathcal{F}$ is relatively compact for weak convergence, i.e. every sequence
$(\mu_n)\subset \mathcal{F}$ admits a weakly convergent subsequence.

More precisely, there exists a probability measure $\mu$ on the Borel $\sigma$-algebra of the completion
$\overline{X}$ of $X$ such that, up to a subsequence,
\[
\mu_{n_k} \rightharpoonup \mu \quad \text{in } \mathcal{P}(\overline{X}).
\]
\end{theorem}

\begin{proof}
Let $(\mu_n)\subset\mathcal{F}$ be an arbitrary sequence.
By tightness, for every $\varepsilon>0$ there exists a compact set $K_\varepsilon\subset X$ such that
\[
\mu_n(K_\varepsilon)\ge 1-\varepsilon
\qquad\text{for all } n.
\]

Choose a sequence $\varepsilon_m\downarrow 0$ and corresponding compact sets $K_m:=K_{\varepsilon_m}$.
Since $K_m$ is compact in the metric space $X$, it is totally bounded.
Hence, for each $m$ and each $\delta>0$, $K_m$ can be covered by finitely many balls of radius $\delta$.

Let $\overline{X}$ denote the completion of $X$.
Each $K_m$ is relatively compact in $\overline{X}$, hence its closure $\overline{K_m}$ is compact in $\overline{X}$.
Moreover,
\[
\mu_n(\overline{K_m}) \ge \mu_n(K_m) \ge 1-\varepsilon_m.
\]

Thus, the sequence $(\mu_n)$ is tight when viewed as a sequence of probability measures on $\overline{X}$.
Since $\overline{X}$ is complete and separable (being the completion of a metric space),
it is a Polish space.

By Prokhorov's theorem on Polish spaces, there exists a subsequence $(\mu_{n_k})$ and a probability measure
$\mu\in\mathcal{P}(\overline{X})$ such that
\[
\mu_{n_k} \rightharpoonup \mu \quad \text{in } \mathcal{P}(\overline{X}).
\]

This proves relative compactness.  
The limit measure $\mu$ need not be supported on $X$, unless $X$ itself is complete.
\end{proof}

\begin{remark}
If $X$ is Polish, then $X=\overline{X}$ and the limit measure belongs to $\mathcal{P}(X)$.
Hence, in this case, tightness is equivalent to relative compactness in $\mathcal{P}(X)$,
which recovers the classical Prokhorov theorem.
\end{remark}


\section{Criteria for tightness}\label{sec:tight-criteria}
A recurring theme in probability theory is the passage to the limit in a sequence of random objects.
For real-valued random variables, compactness arguments are often reduced to showing that the laws do not
``escape to infinity''.  On an infinite-dimensional state space such as the path space
$C([0,\infty);\R^d)$, the situation is more delicate: in addition to controlling the \emph{size} of the paths,
one must also control their \emph{oscillations}.  This is precisely what tightness accomplishes.

\medskip
Recall definition \ref{def:tightness}, we can say that tightness is the probabilistic analogue of precompactness:
it ensures that, along a subsequence, the measures $\Pp_n$ admit a weak limit.
In particular, on Polish spaces, Prokhorov's theorem \ref{thm:prokhorov} shows that tightness is the key compactness
condition for weak convergence of probability measures.

\medskip
In the present setting, we are interested in probability measures on the space of continuous paths
\[
E := C([0,\infty);\R^d),
\]
which typically arise as laws of stochastic processes (e.g.\ approximations of diffusions, rescaled random walks,
numerical schemes, interacting particle systems).  Hence, we need a criterion that allows us to check tightness
directly in terms of path properties of the underlying processes.

\medskip
The guiding idea comes from deterministic compactness.  A set of continuous functions is relatively compact
(with respect to uniform convergence on compact time intervals) if and only if it is \emph{uniformly bounded}
and \emph{equicontinuous}.  This is the content of the Arzel\`a--Ascoli theorem.

\begin{proposition}[Arzel\`a--Ascoli]\label{prop:AA_C0infty}
Let $E:=C([0,\infty);\R^d)$ and let $A\subset E$.
Then $\overline{A}$ is compact in $E$ (with the topology of uniform convergence on compact
intervals) if and only if the following two conditions hold:
\begin{enumerate}[label=\textnormal{(\roman*)}, leftmargin=3.2em]
\item (\emph{Uniform boundedness on compacts}) For every $T>0$,
\[
\sup_{\omega\in A}\|\omega\|_{[0,T]} < \infty,
\qquad\text{where }\ \|\omega\|_{[0,T]}:=\sup_{t\in[0,T]}|\omega(t)|.
\]
\item (\emph{Equicontinuity on compacts}) For every $T>0$,
\[
\lim_{\delta\downarrow 0}\ \sup_{\omega\in A} w_T(\omega,\delta)=0,
\]
with $w_T$ as in Definition~\ref{def:modulus_continuity}.
\end{enumerate}
\end{proposition}

The modulus of continuity defined below provides a convenient way to quantify equicontinuity: small values of $w_T(\omega,\delta)$ for small $\delta$
mean that the path $\omega$ has no rapid oscillations on $[0,T]$.

\begin{definition}[Modulus of continuity]\label{def:modulus_continuity}
Let $E:=C([0,\infty);\R^d)$. For $T>0$ and $\delta>0$ define
\[
w_T(\omega,\delta)
:=\sup\bigl\{|\omega(t)-\omega(s)|:\ s,t\in[0,T],\ |t-s|<\delta\bigr\},
\qquad \omega\in E.
\]
\end{definition}

\medskip
The tightness criterion below can be viewed as a \emph{probabilistic Arzel\`a--Ascoli theorem}.
Condition \emph{(i)} prevents trajectories from becoming arbitrarily large on compact time horizons,
while condition \emph{(ii)} rules out wild oscillations by requiring a uniform (in $n$) control of
the modulus of continuity in probability.  Together these two bounds ensure that, with high probability,
the paths lie in a compact subset of $C([0,\infty);\R^d)$.
This makes the criterion a practical tool for proving weak convergence of processes and is the standard
starting point for functional limit theorems.

\begin{theorem}[Tightness criterion based on Arzel\`a--Ascoli]\label{thm:tightness_C0infty}
Let $E:=C([0,\infty);\R^d)$ and equip $E$ with the metric
\[
d(\omega,\tilde\omega)
:=\sum_{m=1}^\infty 2^{-m}\Bigl(\|\omega-\tilde\omega\|_{[0,m]}\wedge 1\Bigr),
\qquad
\|\omega\|_{[0,T]}:=\sup_{t\in[0,T]}|\omega(t)|.
\]
Then $(E,d)$ is Polish and $\mathcal{B}(E)$ coincides with the canonical $\sigma$-algebra
generated by the coordinate maps (equivalently, by cylinder sets).

Let $(\Pp_n)_{n\in\N}\subset\mathcal{P}(E)$ be a sequence of Borel probability measures.
Then $(\Pp_n)$ is tight in $\mathcal{P}(E)$ if and only if the following two conditions hold:
\begin{enumerate}[label=\textnormal{(\roman*)}, leftmargin=3.2em]
\item (\emph{Tightness of sup-norms on compacts}) For every $T>0$,
\[
\lim_{a\to\infty}\ \sup_{n\in\N}\ \Pp_n\bigl(\|\omega\|_{[0,T]}\ge a\bigr)=0.
\]
\item (\emph{Stochastic equicontinuity on compacts}) For every $T>0$ and every $\varepsilon>0$,
\[
\lim_{\delta\downarrow 0}\ \sup_{n\in\N}\ \Pp_n\bigl(w_T(\omega,\delta)\ge \varepsilon\bigr)=0,
\]
where $w_T$ is the modulus of continuity from Definition~\ref{def:modulus_continuity}.
\end{enumerate}
\end{theorem}

%------------------------------------------------
% Proof of the tightness criterion on C([0,\infty);\R^d)
% (probabilistic translation of Arzelà--Ascoli)
% Based on Section 23 (Theorem 55/56) in the attached notes.
%------------------------------------------------

\begin{proof}[Proof]
We follow the Arzel\`a--Ascoli characterization of compact subsets of
$(C([0,\infty);\R^d),d)$ and translate it into probabilistic language
\medskip
\noindent\textbf{Notation.}
For $T>0$ write $\|\omega\|_{[0,T]}:=\sup_{t\in[0,T]}|\omega(t)|$ and
\[
w_T(\omega,\delta):=\sup\bigl\{|\omega(t)-\omega(s)|:\ s,t\in[0,T],\ |t-s|<\delta\bigr\}.
\]
Recall that a sequence $(\Pp_n)$ is \emph{(uniformly) tight} on $(E,d)$ if for every $\gamma>0$
there exists a compact set $K\subset E$ such that
\[
\inf_{n\in\N}\Pp_n(K)\ge 1-\gamma.
\]

%------------------------------------------------
\medskip
\noindent\textbf{Step 1: ``$\Rightarrow$'' (tightness $\Rightarrow$ the two bounds).}
Assume $(\Pp_n)$ is tight. Fix $\gamma>0$. Then there exists a compact set $K\subset E$ such that
$\Pp_n(K)\ge 1-\gamma$ for all $n$.

\smallskip
\emph{(i) Supremum control.}
Fix $T>0$. Since $K$ is compact, Arzel\`a--Ascoli implies that $K$ is uniformly bounded on $[0,T]$:
there exists $\lambda=\lambda(T)>0$ such that
\[
\|\omega\|_{[0,T]}\le \lambda\qquad\forall\,\omega\in K.
\]
Hence for all $n$,
\[
\Pp_n\bigl(\|\omega\|_{[0,T]}>\lambda\bigr)\le \Pp_n(K^c)\le \gamma.
\]
Letting $\lambda\to\infty$ yields
\[
\lim_{\lambda\to\infty}\ \sup_{n\in\N}\ \Pp_n\bigl(\|\omega\|_{[0,T]}\ge \lambda\bigr)=0.
\]

\smallskip
\emph{(ii) Modulus-of-continuity control.}
Fix $T>0$ and $\varepsilon>0$. Again by Arzel\`a--Ascoli, $K$ is equicontinuous on $[0,T]$:
there exists $\delta_0=\delta_0(T,\varepsilon)>0$ such that for all $\delta\in(0,\delta_0)$,
\[
w_T(\omega,\delta)<\varepsilon\qquad\forall\,\omega\in K.
\]
Thus for all $n$ and all $\delta\in(0,\delta_0)$,
\[
\Pp_n\bigl(w_T(\omega,\delta)\ge\varepsilon\bigr)\le \Pp_n(K^c)\le \gamma.
\]
Letting $\delta\downarrow 0$ (and then $\gamma\downarrow 0$) gives
\[
\lim_{\delta\downarrow 0}\ \sup_{n\in\N}\ \Pp_n\bigl(w_T(\omega,\delta)\ge\varepsilon\bigr)=0.
\]
This proves the two conditions.

%------------------------------------------------
\medskip
\noindent\textbf{Step 2: ``$\Leftarrow$'' (the two bounds $\Rightarrow$ tightness).}
Assume now that the following hold:
\begin{enumerate}[label=\textnormal{(\roman*)}, leftmargin=3.2em]
\item for every $T>0$, $\displaystyle \lim_{\lambda\to\infty}\sup_{n}\Pp_n(\|\omega\|_{[0,T]}\ge\lambda)=0$,
\item for every $T>0$ and $\varepsilon>0$, $\displaystyle \lim_{\delta\downarrow 0}\sup_n\Pp_n(w_T(\omega,\delta)\ge\varepsilon)=0$.
\end{enumerate}
Fix $\gamma>0$. We construct a compact set $K\subset E$ such that $\inf_n\Pp_n(K)\ge 1-\gamma$.

\smallskip
\emph{(a) Compact sets on a fixed horizon $[0,T]$.}
Fix $T\in\N$.
From (i) choose $\lambda_T>0$ such that
\[
\sup_{n\in\N}\Pp_n\bigl(\|\omega\|_{[0,T]}>\lambda_T\bigr)\le \frac{\gamma}{2^{T+1}}.
\]
Next, for each $k\in\N$ apply (ii) with $\varepsilon=1/k$ to find $\delta_{T,k}>0$ such that
\[
\sup_{n\in\N}\Pp_n\bigl(w_T(\omega,\delta_{T,k})>1/k\bigr)\le \frac{\gamma}{2^{T+k+1}}.
\]
Define
\[
A_T
:=\Bigl\{\omega\in E:\ \|\omega\|_{[0,T]}\le \lambda_T,\ 
w_T(\omega,\delta_{T,k})\le \frac1k\ \text{for all }k\in\N\Bigr\}.
\]
Then by the union bound, for all $n\in\N$,
\begin{align*}
\Pp_n(A_T)
&\ge 1-\Pp_n\bigl(\|\omega\|_{[0,T]}>\lambda_T\bigr)
   -\sum_{k=1}^\infty \Pp_n\Bigl(w_T(\omega,\delta_{T,k})>\frac1k\Bigr)\\
&\ge 1-\frac{\gamma}{2^{T+1}}-\sum_{k=1}^\infty \frac{\gamma}{2^{T+k+1}}
= 1-\frac{\gamma}{2^T}.
\end{align*}

\smallskip
\emph{(b) A global compact set.}
Set
\[
K:=\bigcap_{T=1}^\infty A_T.
\]
By construction, $K$ is closed. Moreover, for each fixed $T$, the set $K$ is uniformly bounded
on $[0,T]$ (by $\lambda_T$) and equicontinuous on $[0,T]$ (because
$w_T(\omega,\delta_{T,k})\le 1/k$ for all $k$).
Hence Arzel\`a--Ascoli implies that $K$ is compact in $(E,d)$.

Finally, using again the union bound,
\[
\Pp_n(K)
=1-\Pp_n\Bigl(\bigcup_{T=1}^\infty A_T^c\Bigr)
\ge 1-\sum_{T=1}^\infty \Pp_n(A_T^c)
\ge 1-\sum_{T=1}^\infty \frac{\gamma}{2^T}
=1-\gamma,
\]
for all $n\in\N$. This is exactly tightness.

\medskip
Combining Step~1 and Step~2 proves the theorem.
\end{proof}

\section{Tightness criteria in stochastic processes}
The tightness criterion in Theorem~\ref{thm:tightness_C0infty} is a statement about probability
measures on the path space $E=C([0,\infty);\R^d)$.  In applications, however, we rarely start from
an abstract sequence $(\Pp_n)_{n\in\N}\subset\mathcal{P}(E)$.  Instead, we are typically given a
sequence of \emph{stochastic processes} $(X^n)_{n\in\N}$ defined on (possibly different)
probability spaces, and we are interested in the existence of subsequences converging in
distribution to some limiting process $X$.

The bridge between these two viewpoints is the \emph{law} (distribution) of a process:
whenever $X^n$ has continuous sample paths, its law
\[
\Pp_n:=\mathcal{L}(X^n)\in\mathcal{P}\bigl(C([0,\infty);\R^d)\bigr)
\]
is a well-defined probability measure on the path space.  Thus, to prove convergence in
distribution of processes, it suffices to show tightness of the associated laws $(\Pp_n)$ and then
identify subsequential limits.  Theorem~\ref{thm:tightness_C0infty} reduces this tightness problem
to two intuitive requirements: (i) the trajectories do not escape to infinity on finite time
horizons, and (ii) the trajectories do not oscillate too fast (equicontinuity on compacts).

At this point one encounters a practical issue: conditions \emph{(i)} and \emph{(ii)} in
Theorem~\ref{thm:tightness_C0infty} are formulated in terms of the sup--norm and the modulus of
continuity of the entire sample path $\omega$.  For a given process $X^n$, these are often hard to
estimate directly.  In contrast, one can frequently obtain sharp bounds on the \emph{increments}
$X^n_t-X^n_s$ via moment estimates (e.g.\ from martingale inequalities,
Gaussian estimates, or independence).  The Kolmogorov criterion for tightness is precisely the
tool that translates such increment bounds into the modulus-of-continuity estimate required by
Arzel\`a--Ascoli: moment control of increments implies that, with high probability, paths are
uniformly H\"older continuous on compacts, hence equicontinuous, and therefore tight.

In summary, the change of scope from measure to processes is a
natural viewpoint in probability theory.  It allows us to replace abstract compactness arguments
on $\mathcal{P}(E)$ by concrete, checkable conditions on the processes themselves, ultimately
leading to a practical tightness criterion (Kolmogorov--Chentsov/Kolmogorov tightness criterion)
used throughout weak convergence of stochastic processes.


In order to apply the tightness criterion of Theorem~\ref{thm:tightness_C0infty} to a sequence of
processes $(X^n)_{n\in\N}$, we recall two standard notions identifying processes which coincide ``up to null sets''.

\begin{definition}[Modification]\label{def:modification}
Let $(\Omega,\mathcal{A},\Pp)$ be a probability space and let
$X=(X_t)_{t\ge 0}$ and $Y=(Y_t)_{t\ge 0}$ be $\R^d$-valued stochastic processes.
We say that $X$ and $Y$ are \emph{modifications} (or \emph{versions}) of each other if
for every $t\ge 0$,
\[
\Pp(X_t = Y_t)=1 .
\]
Equivalently, $X_t=Y_t$ $\Pp$-a.s.\ for each fixed $t$.
\end{definition}

\begin{definition}[Indistinguishability]\label{def:indistinguishable}
Let $X=(X_t)_{t\ge 0}$ and $Y=(Y_t)_{t\ge 0}$ be $\R^d$-valued stochastic processes on
$(\Omega,\mathcal{A},\Pp)$.
We say that $X$ and $Y$ are \emph{indistinguishable} if
\[
\Pp\bigl( X_t = Y_t\ \text{for all }t\ge 0 \bigr)=1 .
\]
Equivalently, there exists a null set $N\in\mathcal{A}$ such that for all $\omega\in\Omega\setminus N$,
\[
X_t(\omega)=Y_t(\omega)\qquad\forall\,t\ge 0 .
\]
\end{definition}

\begin{remark}\label{rem:modif-vs-indist}
\begin{enumerate}[label=\textnormal{(\roman*)}, leftmargin=3.2em]
\item Indistinguishability implies modification, but the converse is false in general.
\item If the index set is discrete, e.g.\ $t\in\N$ or $t\in\{0,1,\dots,T\}$, then
\emph{modification and indistinguishability are equivalent}.
Indeed, if $\Pp(X_t=Y_t)=1$ for each $t$ in a countable set, then the intersection of the
corresponding full-measure events is still a full-measure event.
\item If $X$ and $Y$ admit continuous sample paths and are modifications of each other, then
they are automatically indistinguishable.
\end{enumerate}
\end{remark}

\medskip
Furthermore, to apply the Kolmogorov criterion, we need to introduce (remember) H\"older continuity.

\begin{definition}[Local H\"older continuity]\label{def:holder_local}
Let $\alpha\in(0,1]$ and $T>0$.
A function $f:[0,T]\to\R^d$ is called \emph{$\alpha$-H\"older continuous on $[0,T]$} if
there exists a constant $C_T<\infty$ such that
\[
|f(t)-f(s)| \le C_T\,|t-s|^\alpha
\qquad\forall\, s,t\in[0,T].
\]
The smallest such constant is denoted by
\[
[f]_{\alpha;[0,T]} := \sup_{\substack{s,t\in[0,T]\\s\neq t}}
\frac{|f(t)-f(s)|}{|t-s|^\alpha}
\in [0,\infty],
\]
and is called the \emph{H\"older seminorm} of $f$ on $[0,T]$.
\end{definition}

\begin{definition}[H\"older continuity]\label{def:holder_global}
A function $f:[0,\infty)\to\R^d$ is called \emph{locally $\alpha$-H\"older continuous} if
for every $T>0$ the restriction $f|_{[0,T]}$ is $\alpha$-H\"older continuous
(i.e.\ $[f]_{\alpha;[0,T]}<\infty$ for all $T>0$).
\end{definition}

\begin{remark}[H\"older continuity implies modulus of continuity estimates]\label{rem:holder_modulus}
If $f$ is $\alpha$-H\"older continuous on $[0,T]$ with constant $C_T$, then for every $\delta>0$,
\[
w_T(f,\delta)\le C_T\,\delta^\alpha,
\]
where $w_T$ is the modulus of continuity from Definition~\ref{def:modulus_continuity}.
In particular, H\"older regularity implies equicontinuity, which is the key compactness
mechanism behind the tightness criterion of Theorem~\ref{thm:tightness_C0infty}.
\end{remark}

\medskip
The criterion is indeed split in two theorems. The first -- Kolmogorov--Chentsov theorem -- will provide sufficient moment conditions on the increments
of a stochastic process guaranteeing the existence of a \emph{continuous} (in fact H\"older continuous)
modification. This is the main tool for producing probability measures on
$C([0,\infty);\R^d)$ starting from a process given only by its finite-dimensional distributions.

\begin{theorem}[Kolmogorov--Chentsov]\label{thm:kolmogorov_chentsov}
Let $(\Omega,\mathcal{A},\Pp)$ be a probability space and let
$X=(X_t)_{t\in[0,T]}$ be an $\R^d$-valued stochastic process.
Assume that there exist constants $p>0$, $\beta>0$ and $C>0$ such that
\begin{equation}\label{eq:KC_assumption}
\E\bigl[|X_t-X_s|^p\bigr]\le C\,|t-s|^{1+\beta}
\qquad\forall\, s,t\in[0,T].
\end{equation}
Then there exists a modification $\widetilde X$ of $X$ such that
for every $\alpha\in\bigl(0,\beta/p\bigr)$ there exists an a.s.\ finite random variable
$K_\alpha$ with
\begin{equation}\label{eq:KC_holder}
|\widetilde X_t(\omega)-\widetilde X_s(\omega)|
\le K_\alpha(\omega)\,|t-s|^\alpha
\qquad \forall\, s,t\in[0,T],\ \text{for a.e.\ }\omega.
\end{equation}
In particular, $\widetilde X$ has (a.s.) continuous sample paths on $[0,T]$.

Moreover, repeating the argument on $[0,m]$ for $m\in\N$ shows that one can obtain a modification
on $[0,\infty)$ whose paths are locally $\alpha$-H\"older continuous for every $\alpha<\beta/p$.
\end{theorem}

\begin{proof}[Proof]
The proof is a quantitative version of the idea
\emph{``moment control of increments $\Rightarrow$ small increments with high probability
$\Rightarrow$ only finitely many bad oscillations $\Rightarrow$ H\"older continuity.''}

\medskip
\noindent\LabelProof{Step 1: Reduction to the scalar case and choice of exponent}
Since $|x| \le \sum_{j=1}^d |x_j|$ and $(a_1+\cdots+a_d)^p \le d^{p-1}(a_1^p+\cdots+a_d^p)$ for $p\ge 1$
(one can treat $p\in(0,1)$ similarly),
it suffices to prove the statement for real-valued processes and the same moment bound
(up to changing the constant $C$).
Fix once and for all
\[
\alpha\in\bigl(0,\beta/p\bigr),
\qquad \text{and choose }\ \theta>0\ \text{such that }\ \alpha=\frac{\beta-\theta}{p}.
\]
Equivalently, $p\alpha=\beta-\theta$ with $\theta\in(0,\beta)$.

\medskip
\noindent\LabelProof{Step 2: Dyadic increments and a key probability estimate}
For each level $m\in\N$ consider the dyadic grid
\[
D_m:=\Bigl\{\frac{kT}{2^m}: k=0,1,\dots,2^m\Bigr\}.
\]
Define the maximal dyadic increment on this grid:
\[
\Delta_m := \max_{0\le k\le 2^m-1}\Bigl|X_{\frac{(k+1)T}{2^m}}-X_{\frac{kT}{2^m}}\Bigr|.
\]
We will show that $\Delta_m$ decays roughly like $2^{-m\alpha}$.

Set $\lambda_m := 2^{-m\alpha}$.
Using Markov's inequality and the union bound,
\begin{align*}
\Pp\bigl(\Delta_m > \lambda_m\bigr)
&\le \sum_{k=0}^{2^m-1}
\Pp\Bigl(\Bigl|X_{\frac{(k+1)T}{2^m}}-X_{\frac{kT}{2^m}}\Bigr|>\lambda_m\Bigr)\\
&\le \sum_{k=0}^{2^m-1}
\frac{\E\Bigl[\bigl|X_{\frac{(k+1)T}{2^m}}-X_{\frac{kT}{2^m}}\bigr|^p\Bigr]}{\lambda_m^p}.
\end{align*}
By the increment moment bound \eqref{eq:KC_assumption},
\[
\E\Bigl[\Bigl|X_{\frac{(k+1)T}{2^m}}-X_{\frac{kT}{2^m}}\Bigr|^p\Bigr]
\le C\Bigl(\frac{T}{2^m}\Bigr)^{1+\beta}.
\]
Hence
\begin{align*}
\Pp(\Delta_m > \lambda_m)
&\le 2^m\,
\frac{C\,(T/2^m)^{1+\beta}}{(2^{-m\alpha})^p}
= C\,T^{1+\beta}\,2^m\,2^{-m(1+\beta)}\,2^{mp\alpha}\\
&= C\,T^{1+\beta}\,2^{-m\bigl(\beta-p\alpha\bigr)}
= C\,T^{1+\beta}\,2^{-m\theta}.
\end{align*}
Since $\theta>0$, the series $\sum_{m=1}^\infty 2^{-m\theta}$ converges. Therefore,
\begin{equation}\label{eq:BC_summable}
\sum_{m=1}^\infty \Pp(\Delta_m > 2^{-m\alpha}) < \infty.
\end{equation}

\medskip
\noindent\LabelProof{Step 3: Borel--Cantelli $\Rightarrow$ dyadic increments are eventually small}
By \eqref{eq:BC_summable} and the Borel--Cantelli lemma,
\[
\Pp\bigl(\Delta_m > 2^{-m\alpha}\ \text{i.o.}\bigr)=0.
\]
Thus there exists a full-measure event $\Omega_0$ such that for every $\omega\in\Omega_0$
there exists $m_0(\omega)$ with
\begin{equation}\label{eq:dyadic_control}
\Delta_m(\omega)\le 2^{-m\alpha}\qquad\forall m\ge m_0(\omega).
\end{equation}

Intuition: for large $m$ the mesh size $T/2^m$ is tiny, and the moment condition forces
\emph{most} increments to be tiny; Borel--Cantelli tells us that eventually
\emph{all} increments on the dyadic grid are tiny.

\medskip
\noindent\LabelProof{Step 4: H\"older control on dyadic rationals}
Fix $\omega\in\Omega_0$.
Let $s<t$ be points in the union of dyadic grids
\[
D:=\bigcup_{m=1}^\infty D_m,
\]
and choose $m$ so large that $s,t\in D_m$.
Write $s=\frac{iT}{2^m}$ and $t=\frac{jT}{2^m}$ with $i<j$.
Then by telescoping,
\[
|X_t(\omega)-X_s(\omega)|
\le \sum_{\ell=i}^{j-1}\Bigl|X_{\frac{(\ell+1)T}{2^m}}(\omega)-X_{\frac{\ell T}{2^m}}(\omega)\Bigr|
\le (j-i)\,\Delta_m(\omega).
\]
Using \eqref{eq:dyadic_control} (for large enough $m$),
\[
|X_t(\omega)-X_s(\omega)|
\le (j-i)\,2^{-m\alpha}
= \Bigl(\frac{(j-i)T}{2^m}\Bigr)^\alpha \,T^{1-\alpha}\,(j-i)^{1-\alpha}.
\]
Since $(j-i)T/2^m = t-s$ and $(j-i)^{1-\alpha}\le \bigl(2^m(t-s)/T\bigr)^{1-\alpha}$, we obtain
(after a short rearrangement) the bound
\[
|X_t(\omega)-X_s(\omega)|
\le C_\alpha(\omega)\,|t-s|^\alpha
\qquad\forall\, s,t\in D,
\]
for some finite constant $C_\alpha(\omega)$ depending on $\omega$ and $\alpha$.
(One convenient choice is
$C_\alpha(\omega):=\sup_{m\ge m_0(\omega)} 2^{m\alpha}\Delta_m(\omega)$, which is finite by
\eqref{eq:dyadic_control}.)

\medskip
\noindent\LabelProof{Step 5: Construction of a continuous modification}
For $\omega\in\Omega_0$, the map $t\mapsto X_t(\omega)$ is now H\"older continuous on the dense set $D$.
Hence it has a unique continuous extension to $[0,T]$, which we denote by
\[
\widetilde X_t(\omega):=\lim_{\substack{r\to t\\ r\in D}} X_r(\omega),
\qquad t\in[0,T],\ \omega\in\Omega_0.
\]
On the null set $\Omega\setminus\Omega_0$ define $\widetilde X(\omega)$ arbitrarily (e.g.\ $0$).

Then $\widetilde X(\omega)$ is continuous for every $\omega\in\Omega_0$ by construction,
and the H\"older estimate \eqref{eq:KC_holder} extends from $D$ to all $s,t\in[0,T]$
by density and continuity.

\medskip
\noindent\LabelProof{Step 6: $\widetilde X$ is a modification of $X$}
Fix any $t\in[0,T]$.
Choose a sequence $(r_n)_{n\in\N}\subset D$ with $r_n\to t$.
By definition of $\widetilde X_t$ and by \eqref{eq:KC_assumption},
\[
\E\bigl[|X_{r_n}-X_t|^p\bigr]\le C|r_n-t|^{1+\beta}\to 0,
\]
so $X_{r_n}\to X_t$ in $L^p$, hence (along a subsequence) almost surely.
But on $\Omega_0$ we also have $X_{r_n}\to \widetilde X_t$ by definition.
Therefore $X_t=\widetilde X_t$ almost surely, i.e.\ $\widetilde X$ is a modification of $X$.

\medskip
This proves the theorem on $[0,T]$. For $[0,\infty)$ one applies the argument on each $[0,m]$,
$m\in\N$, and intersects the full-measure events to obtain local H\"older continuity.
\end{proof}


The second theorem -- Kolmogorov tightness criterion -- uses the Kolmogorov--Chentsov theorem to give
sufficient conditions for tightness of a sequence of processes with continuous sample paths.

%------------------------------------------------
% Kolmogorov criterion for tightness on C([0,T];R^d)
% (stepwise + intuitive proof, in the style used above)
%------------------------------------------------

\begin{theorem}[Kolmogorov criterion for tightness]\label{thm:kolmogorov_tightness}
Let $T>0$ and let $(X^n)_{n\in\N}$ be a sequence of $\R^d$-valued stochastic processes on $[0,T]$
(with measurable sample paths).
Assume that:
\begin{enumerate}[label=\textnormal{(\roman*)}, leftmargin=3.2em]
\item \emph{Tight initial laws:} the family $\{\mathcal{L}(X_0^n)\}_{n\in\N}$ is tight in $\R^d$.
\item \emph{Uniform increment moment bound:} there exist constants $p>0$, $\beta>0$ and $C>0$ such that
\begin{equation}\label{eq:kolmo_moment_bound}
\E\bigl[|X_t^n-X_s^n|^p\bigr]\le C\,|t-s|^{1+\beta}
\qquad\forall\, s,t\in[0,T],\ \forall\,n\in\N.
\end{equation}
\end{enumerate}
Then the sequence of laws $\Pp_n:=\mathcal{L}(X^n)$ is tight in $\mathcal{P}(C([0,T];\R^d))$.

In particular, if each $X^n$ has a continuous modification (e.g.\ by Kolmogorov--Chentsov),
then we may assume $X^n\in C([0,T];\R^d)$ and tightness holds in the path space.
\end{theorem}

\begin{proof}[Proof (following the Arzel\`a--Ascoli tightness criterion)]
We verify the two conditions of Theorem~\ref{thm:tightness_C0infty} on the finite horizon $[0,T]$,
namely:
\[
\text{(A) tightness of }\sup_{t\in[0,T]}|X_t^n|
\qquad\text{and}\qquad
\text{(B) stochastic equicontinuity via }w_T(X^n,\delta).
\]
Throughout the proof, constants may depend on $T,p,\beta,C$ but never on $n$.

\medskip
\noindent\LabelProof{Step 1: A dyadic increment estimate (uniform in $n$)}
Fix $\alpha\in(0,\beta/p)$ and set $\theta:=\beta-p\alpha>0$.
For $m\in\N$ consider the dyadic grid $t_k:=kT2^{-m}$, $k=0,\dots,2^m$, and define
\[
\Delta_m^n:=\max_{0\le k\le 2^m-1}\bigl|X^n_{t_{k+1}}-X^n_{t_k}\bigr|.
\]
Then for every $\lambda>0$ (by union bound and Markov's inequality),
\begin{align*}
\Pp\bigl(\Delta_m^n>\lambda\bigr)
&\le \sum_{k=0}^{2^m-1} \Pp\bigl(|X^n_{t_{k+1}}-X^n_{t_k}|>\lambda\bigr)\\
&\le \sum_{k=0}^{2^m-1} \frac{\E\bigl[|X^n_{t_{k+1}}-X^n_{t_k}|^p\bigr]}{\lambda^p}
\le 2^m\frac{C (T2^{-m})^{1+\beta}}{\lambda^p}.
\end{align*}
Choosing $\lambda=2^{-m\alpha}$ yields
\begin{equation}\label{eq:dyadic_tail_bound}
\sup_{n\in\N}\Pp\bigl(\Delta_m^n>2^{-m\alpha}\bigr)
\le C\,T^{1+\beta}\,2^{-m(\beta-p\alpha)}
= C\,T^{1+\beta}\,2^{-m\theta}.
\end{equation}
This is the basic estimate: \emph{large oscillations on a very fine grid are uniformly unlikely.}

\medskip
\noindent\LabelProof{Step 2: Stochastic equicontinuity (condition (B))}
Fix $\varepsilon>0$.
Let $m\in\N$ and choose $\delta_m:=T2^{-m}$.
If $|t-s|<\delta_m$, then $s$ and $t$ lie in at most two neighbouring dyadic intervals of level $m$,
hence one has the deterministic bound
\begin{equation}\label{eq:modulus_vs_dyadic}
w_T(X^n,\delta_m)\le 2\,\Delta_m^n.
\end{equation}
Therefore, using \eqref{eq:dyadic_tail_bound},
\[
\sup_{n\in\N}\Pp\bigl(w_T(X^n,\delta_m)\ge\varepsilon\bigr)
\le \sup_{n\in\N}\Pp\Bigl(\Delta_m^n\ge \frac{\varepsilon}{2}\Bigr).
\]
Now pick $m$ so large that $2^{-m\alpha}\le \varepsilon/2$; then
\[
\Pp\Bigl(\Delta_m^n\ge \frac{\varepsilon}{2}\Bigr)
\le \Pp\bigl(\Delta_m^n\ge 2^{-m\alpha}\bigr).
\]
Hence by \eqref{eq:dyadic_tail_bound},
\[
\sup_{n\in\N}\Pp\bigl(w_T(X^n,\delta_m)\ge\varepsilon\bigr)
\le C\,T^{1+\beta}\,2^{-m\theta}.
\]
Letting $m\to\infty$ (equivalently $\delta_m\downarrow 0$) gives
\[
\lim_{\delta\downarrow 0}\ \sup_{n\in\N}\Pp\bigl(w_T(X^n,\delta)\ge\varepsilon\bigr)=0,
\]
which is exactly stochastic equicontinuity.

\medskip
\noindent\LabelProof{Step 3: Tightness of the sup--norm (condition (A))}
Fix $\eta>0$. We will show that there exists $M>0$ such that
\[
\sup_{n\in\N}\Pp\Bigl(\sup_{t\in[0,T]}|X_t^n|>M\Bigr)\le \eta.
\]

\smallskip
\emph{Step 3a: Control the initial value.}
By tightness of $\{\mathcal{L}(X_0^n)\}$, choose $R>0$ such that
\begin{equation}\label{eq:init_tight}
\sup_{n\in\N}\Pp(|X_0^n|>R)\le \eta/2.
\end{equation}

\smallskip
\emph{Step 3b: Control the increments on a coarse grid.}
Pick $m\in\N$ large enough so that
\[
\sup_{n\in\N}\Pp\bigl(w_T(X^n,\delta_m)>1\bigr)\le \eta/4,
\qquad \delta_m:=T2^{-m},
\]
which is possible by Step~2 with $\varepsilon=1$.
Let $t_k:=k\delta_m$, $k=0,\dots,2^m$. On the event $\{w_T(X^n,\delta_m)\le 1\}$ we have
\[
\sup_{t\in[0,T]}|X_t^n|
\le \max_{0\le k\le 2^m}|X_{t_k}^n| + 1.
\]
Thus it remains to control the maximum over the \emph{finite} set of grid values.

Using the increment bound \eqref{eq:kolmo_moment_bound} and Markov's inequality,
for any $a>0$ and each $k$,
\[
\Pp\bigl(|X_{t_k}^n-X_0^n|>a\bigr)
\le \frac{\E\bigl[|X_{t_k}^n-X_0^n|^p\bigr]}{a^p}
\le \frac{C\,t_k^{1+\beta}}{a^p}
\le \frac{C\,T^{1+\beta}}{a^p}.
\]
By union bound over $k=0,\dots,2^m$ we obtain
\[
\Pp\Bigl(\max_{0\le k\le 2^m}|X_{t_k}^n-X_0^n|>a\Bigr)
\le 2^m\,\frac{C\,T^{1+\beta}}{a^p}.
\]
Choose $a=a(\eta,m)$ so large that
\begin{equation}\label{eq:grid_tight}
2^m\,\frac{C\,T^{1+\beta}}{a^p}\le \eta/4.
\end{equation}

\smallskip
\emph{Step 3c: Combine the estimates.}
Set $M:=R+a+1$. Then
\begin{align*}
\Pp\Bigl(\sup_{t\in[0,T]}|X_t^n|>M\Bigr)
&\le \Pp(|X_0^n|>R)
 +\Pp\Bigl(\max_{k}|X_{t_k}^n-X_0^n|>a\Bigr)
 +\Pp\bigl(w_T(X^n,\delta_m)>1\bigr)\\
&\le \eta/2+\eta/4+\eta/4=\eta,
\end{align*}
uniformly in $n$, by \eqref{eq:init_tight}, \eqref{eq:grid_tight} and the choice of $m$.
This proves (A).

\medskip
\noindent\LabelProof{Step 4: Conclusion (tightness)}
Steps~2 and 3 verify the two conditions of the Arzel\`a--Ascoli tightness criterion
(Theorem~\ref{thm:tightness_C0infty} restricted to $[0,T]$).
Hence the sequence of laws $\Pp_n=\mathcal{L}(X^n)$ is tight in
$\mathcal{P}(C([0,T];\R^d))$.
\end{proof}

\section{Conclusion on weak convergence}

Weak convergence provides a robust and flexible notion of convergence for probability measures: instead of requiring
$\mu_n(A)\to\mu(A)$ for all measurable sets (too rigid in practice), it asks that expectations of \emph{bounded continuous observables}
stabilize, i.e.\ $\int f\,d\mu_n\to\int f\,d\mu$ for all $f\in C_b(E)$.
This viewpoint is both conceptually natural (``distributions are determined by what they predict on observables'')
and technically powerful, because $C_b(E)$ is separating: if two measures agree on all bounded continuous tests,
then they are the same.

A key message of this chapter is that weak convergence can be recognized in several equivalent ways.
The Portmanteau theorem gives a dictionary between the test-function formulation and setwise criteria:
weak convergence is equivalent to one-sided bounds on probabilities of closed/open sets and to convergence on
$\mu$-continuity sets (where boundaries carry no mass).
The Prokhorov metric then turns this topology into a geometric, quantitative notion of distance between measures,
so that one can speak of Cauchy sequences and stability of limits in a metric framework.

Finally, on non-compact spaces, weak convergence requires a mechanism preventing ``escape of mass to infinity''.
Tightness is exactly this mechanism, and Prokhorov’s theorem identifies it as the precise compactness criterion:
on Polish spaces, tightness $\Longleftrightarrow$ relative weak compactness.
In practice, this means that to prove convergence in law one often proceeds in two steps:
(i) prove tightness (compactness), and (ii) identify the limit by checking convergence on a convergence-determining class.

Weak convergence is therefore indispensable across analysis and probability:
it is the language behind convergence in distribution (CLT and functional limit theorems),
markov chain approximations in stochastic control, non-parametric stochastics for continuous processes and many other applications.
In all these settings, weak convergence captures the correct ``macroscopic'' notion of convergence:
what matters is that all bounded continuous observables have asymptotically the same expectation.







